
% Requires in the main manuscript preamble:
% \usepackage{booktabs}
% \usepackage{pdflscape}
% \usepackage{longtable}
% \usepackage{array}

\begin{landscape}

% If this is the first supplementary table, place these two lines
% immediately before the table (or set them once at the start of the supplement):
% \setcounter{table}{0}
% \renewcommand{\thetable}{S\arabic{table}}

\small
\setlength{\tabcolsep}{4pt}
\renewcommand{\arraystretch}{1.15}
\setlength{\LTcapwidth}{0.95\linewidth}

\begin{longtable}{
    >{\raggedright\arraybackslash}p{3.15cm}
    >{\centering\arraybackslash}p{1.35cm}
    >{\raggedright\arraybackslash}p{3.15cm}
    >{\raggedright\arraybackslash}p{6.15cm}
    >{\raggedright\arraybackslash}p{6.0cm}
}

\caption{Sensitivity and robustness analyses of age modification in the HOMA-IR–hs-CRP association, NHANES 2015–March 2020
}
\label{tab:S1}\\

\toprule
\textbf{Analysis} &
\textbf{$n$} &
\textbf{Interaction / test} &
\textbf{Age-specific HOMA-IR coefficients} &
\textbf{Interpretation} \\
\midrule
\endfirsthead

\multicolumn{5}{l}{\textit{Supplementary Table \thetable\ continued}}\\[2pt]
\toprule
\textbf{Analysis} &
\textbf{$n$} &
\textbf{Interaction / test} &
\textbf{Age-specific HOMA-IR coefficients} &
\textbf{Interpretation} \\
\midrule
\endhead

\midrule
\multicolumn{5}{r}{\textit{Continued on next page}}\\
\endfoot

\bottomrule
\endlastfoot

\textbf{Primary model} &
3,079 &
$F(2,40)=5.69$, $p=0.00671$ &
20--39: 0.0266 (0.0038, 0.0495)\newline
40--54: 0.0511 (0.0132, 0.0890)\newline
55+: $-0.0189$ ($-0.0513$, 0.0135) &
Reference specification. \\

\addlinespace

Raw hs-CRP outcome &
3,079 &
$F(2,40)=3.96$, $p=0.0269$ &
20--39: 0.127 (0.018, 0.235)\newline
40--54: 0.160 ($-0.056$, 0.377)\newline
55+: $-0.109$ ($-0.282$, 0.064) &
Age interaction retained on the original outcome scale. \\

\addlinespace

Log-transformed HOMA-IR &
3,079 &
$F(2,40)=1.86$, $p=0.169$ &
20--39: 0.127 ($-0.0068$, 0.261)\newline
40--54: 0.172 ($-0.0059$, 0.350)\newline
55+: $-0.0046$ ($-0.153$, 0.144) &
Directional pattern remained, but formal age-interaction evidence was attenuated. \\

\addlinespace

Raw hs-CRP + log HOMA-IR &
3,079 &
$F(2,40)=1.43$, $p=0.251$ &
20--39: 0.475 ($-0.146$, 1.096)\newline
40--54: 0.338 ($-0.621$, 1.296)\newline
55+: $-0.165$ ($-0.937$, 0.607) &
Formal interaction not detected. \\

\addlinespace

No continuous HbA1c adjustment &
3,079 &
$F(2,40)=5.57$, $p=0.00732$ &
20--39: 0.0284\newline
40--54: 0.0521\newline
55+: $-0.0171$ &
Essentially unchanged from the primary model. \\

\addlinespace

BMI substituted for waist circumference\textsuperscript{a} &
3,072 &
$F(2,40)=4.91$, $p=0.0124$ &
20--39: 0.0304 (0.0058, 0.0550)\newline
40--54: 0.0601 (0.0234, 0.0968)\newline
55+: $-0.0085$ ($-0.0468$, 0.0298) &
Age pattern retained under alternative adiposity adjustment. \\

\addlinespace

hs-CRP $\leq10$ mg/L, log outcome &
2,895 &
$F(2,40)=2.58$, $p=0.0883$ &
20--39: 0.0249 (0.0054, 0.0443)\newline
40--54: 0.0478 (0.0118, 0.0838)\newline
55+: $-0.0042$ ($-0.0367$, 0.0283) &
Age-specific pattern remained, but formal interaction evidence weakened. \\

\addlinespace

hs-CRP $\leq10$ mg/L, raw outcome &
2,895 &
$F(2,40)=4.03$, $p=0.0255$ &
20--39: 0.0824 (0.0300, 0.1348)\newline
40--54: 0.0905 (0.0108, 0.1702)\newline
55+: $-0.0304$ ($-0.1077$, 0.0469) &
Interaction retained after restricting inflammatory extremes. \\

\addlinespace

HbA1c $<6.0\%$ &
4,018 &
$F(2,40)=5.74$, $p=0.00644$ &
20--39: 0.0185 ($-0.0034$, 0.0404)\newline
40--54: 0.0448 (0.0070, 0.0826)\newline
55+: $-0.0190$ ($-0.0457$, 0.0078) &
Relaxing eligibility did not remove age modification. \\

\addlinespace

HbA1c $<6.5\%$ &
4,492 &
$F(2,40)=5.58$, $p=0.00731$ &
20--39: 0.0120 ($-0.0101$, 0.0341)\newline
40--54: 0.0371 (0.0046, 0.0696)\newline
55+: $-0.0290$ ($-0.0568$, $-0.0011$) &
Age modification persisted with a broader glycemic range. \\

\addlinespace

No HbA1c cutoff\textsuperscript{b} &
4,639 &
$F(2,40)=5.74$, $p=0.00643$ &
20--39: 0.0115 ($-0.0081$, 0.0311)\newline
40--54: 0.0349 (0.0065, 0.0633)\newline
55+: $-0.0212$ ($-0.0456$, 0.0032) &
Simple HbA1c-threshold selection did not explain older attenuation. \\

\addlinespace

Continuous linear age interaction &
3,079 &
$F(1,40)=3.93$, $p=0.0544$ &
Not represented by three age-group coefficients. &
Borderline evidence under a strictly linear age-varying slope. \\

\addlinespace

Continuous quadratic age interaction &
3,079 &
$F(2,40)=3.43$, $p=0.0421$ &
Not represented by three age-group coefficients. &
Evidence of nonlinear age modification; quadratic component alone $p=0.0857$. \\

\addlinespace

Natural cubic spline age interaction &
3,079 &
$F(4,40)=3.09$, $p=0.0262$ &
Not represented by three age-group coefficients. &
Flexible continuous-age model supported age-dependent variation. \\

\addlinespace

Leave-one-survey-cluster-out\textsuperscript{c} &
3,023--3,063 &
79 refits &
Not applicable. &
Interaction $F$ ranged from 3.59 to 7.66 and $p$ from 0.00157 to 0.0369.
No single masked stratum--PSU cluster eliminated the interaction. \\

\end{longtable}

\vspace{-0.5em}
\begin{minipage}{0.97\linewidth}
\footnotesize
\textit{Notes:} Values in parentheses are 95\% confidence intervals.
For models with $\ln$(hs-CRP) as the outcome, coefficients are on the log outcome scale;
raw-outcome coefficients are in hs-CRP units.
All analyses used the NHANES complex survey design with duration-adjusted fasting-subsample weights,
masked variance strata, and masked variance units.
Sensitivity analyses were performed to evaluate robustness and were not treated as additional primary hypothesis tests.

\textsuperscript{a} To isolate the adiposity specification itself, the BMI analysis was restricted
to participants in the primary analytic sample who also had nonmissing BMI.

\textsuperscript{b} Participants with diagnosed diabetes remained excluded; only the HbA1c
eligibility threshold was removed.

\textsuperscript{c} Each refit deleted one unique masked stratum--PSU combination from the locked primary model.
\end{minipage}

\end{landscape}
